\section*{Chapter \ref{ch:mx:fees-rebates-and-venue-economics} --- Fees, Rebates and Venue Economics}

\begin{solution}{exo:mx:fees-rebates-and-venue-economics:1}
BZX: $20.00+0.0030=20.0030$; BYX: $20.00-0.0002=19.9998$. The inverted venue saves 0.32 cents a share, USD 3.20 on 1\,000 shares.
\end{solution}

\begin{solution}{exo:mx:fees-rebates-and-venue-economics:2}
The fee-adjusted asks interleave with gaps $0.30-(-0.02)=0.32$ cents and $1-0.32=0.68$ cents: the effective tick is 0.32 cents.
\end{solution}

\begin{solution}{exo:mx:fees-rebates-and-venue-economics:3}
INV: $0.3\times(0.5-0.1-0.2)=0.06$ cents; MT: $0.1\times(0.5-0.6+0.16)=0.006$ cents.
\end{solution}

\begin{solution}{exo:mx:fees-rebates-and-venue-economics:4}
Solve $0.1\,(-0.1-m)=0.06$: $m=-0.7$, a rebate of 0.70 cents. The venue would pay 0.70 to the maker and receive 0.30 from the taker, losing 0.40 cents on
every share traded: it cannot.
\end{solution}

\begin{solution}{exo:mx:fees-rebates-and-venue-economics:5}
The total fee stays at 0.14 cents. Without a tick the ask moves to 10.012: the taker pays $10.012+0.001=10.013$ as before ($10.01+0.003$), the maker nets
$10.012-0.0004=10.0116$ as before ($10.01+0.0016$). With a one-cent tick the ask stays at 10.01 (taker 10.011, the maker loses 0.2 cents) or moves to 10.02
(taker 10.021): the split is no longer \omterm{def:mx:fees-rebates-and-venue-economics:neutral}{neutral}.
\end{solution}

\begin{solution}{exo:mx:fees-rebates-and-venue-economics:6}
The capture is the half-spread, the same at the same price on both venues. What changes is which trades reach MT's queue: fee-aware takers hit INV first,
so MT fills only when INV's queue has been exhausted, by large or informed orders after which the price moves on. The more takers are fee-aware, the
more MT's fills are these.
\end{solution}

\begin{solution}{exo:mx:fees-rebates-and-venue-economics:7}
Three seeds of two hours: MT's filled share (8.4\%) and adverse selection (0.63 cents) do not change, because INV's take rebate still makes it cheaper to
take and the routing order is the same; the value per posted share falls from 0.006 to $-0.003$ cents, the 0.10 cents of rebate lost on each filled share.
A lower cap removes room for rebates without changing who fills first.
\end{solution}

\begin{solution}{exo:mx:fees-rebates-and-venue-economics:8}
The rebate goes to the broker, and the order stands where fee-aware takers arrive last: it fills less often and, when it fills, more often against flow
that keeps moving the price. Battalio, Corwin and Jennings found a negative relation between \omterm{def:mx:the-limit-order-book:orders}{limit order} execution quality and the rebate level; the
simulation shows fills of 8.4\% against 28.2\% and adverse selection of 0.63 against 0.13 cents. Measure fill rates and mark-outs by venue before claiming
lower costs.
\end{solution}

\begin{solution}{pb:mx:fees-rebates-and-venue-economics:1}
\textbf{1.} BZX: rebate 0.16 cents to add, fee 0.30 to remove; BYX: fee 0.20 to add, rebate 0.02 to remove.
\textbf{2.} 20.0030 and 19.9998.
\textbf{3.} The smallest gap between \omterm{def:mx:fees-rebates-and-venue-economics:adjusted}{fee-adjusted prices} across venues: 0.32 cents.
\textbf{4.} Venues cannot quote between ticks; different fees let them offer sub-tick prices net of fees (Chao, Yao and Ye).
\textbf{5.} See the proposition: $a+f^{take}=R+F$ without a tick; with one, $a$ is rounded up to the grid.
\textbf{6.} Without a tick the ask moves to 10.012 and nothing else changes; with one it stays at 10.01 or moves to 10.02.
\textbf{7.} To queue position: makers compete by joining queues, so the split changes queue lengths.
\textbf{8.} Quotes adjusted and demanders' costs, fees included, did not change; posted spreads narrowed and aggressive orders became more frequent.
\textbf{9.} An efficient price with jumps, \omterm{def:mx:the-limit-order-book:orders}{limit orders} at or behind the best cancelling at a constant rate and when stale, noise and informed \omterm{def:mx:the-limit-order-book:orders}{market orders};
fee-aware takers route by \omterm{def:mx:fees-rebates-and-venue-economics:adjusted}{fee-adjusted price}, fee-blind ones by displayed size.
\textbf{10.} INV 77\% of the volume, 28.2\% filled, 4.6 seconds; MT 23\%, 8.4\%, 8.8 seconds.
\textbf{11.} 0.63 cents on MT against 0.13 on INV: MT fills when INV's queue is gone, which is when large or informed flow arrives.
\textbf{12.} 0.006 cents on MT, 0.057 on INV.
\textbf{13.} A rebate of 0.77 cents, beyond the 0.30-cent take fee: no.
\textbf{14.} MT's queue is taken first by size-based routers: 0.052 cents on MT against 0.028 on INV, and the indifferent make fee on MT is about zero.
\textbf{15.} 0.074 against 0.015 cents.
\textbf{16.} Rule 610T, adopted December 19, 2018: 1\,460 stocks in a USD 0.0010 fee-cap group and a no-rebate group; vacated by the D.C. Circuit on
June 16, 2020.
\textbf{17.} Brokers sending \omterm{def:mx:the-limit-order-book:orders}{limit orders} to the highest-rebate venues, and worse \omterm{def:mx:the-limit-order-book:orders}{limit order} execution quality where rebates are higher.
\textbf{18.} Rebates must stay below the 0.10-cent take fee, so the most a maker-taker venue can pay falls by 0.20 cents.
\textbf{19.} \emph{Named result}: with all takers fee-aware, a posted share fills 8.4\% of the time on the maker-taker venue with 0.63 cents of 30-second
adverse selection, against 28.2\% and 0.13 cents on the inverted venue; a provider would be indifferent only at a 0.77-cent rebate, which the venue cannot
pay; with half the takers routing by size the indifferent make fee is about zero.
\textbf{20.} A place at the back of the consolidated queue, worth what the takers' routing makes it worth.
\end{solution}

\begin{solution}{iq:mx:fees-rebates-and-venue-economics:1}
Maker-taker pays the resting side and charges the taker; inverted does the reverse. To be filled quickly, post on the inverted venue: fee-aware takers
hit it first at equal prices, and its queue is shorter.
\iqlookfor{Fee roles; routing order; queue length.}
\end{solution}

\begin{solution}{iq:mx:fees-rebates-and-venue-economics:2}
Without a tick, yes: quotes absorb the move and \omterm{def:mx:fees-rebates-and-venue-economics:adjusted}{fee-adjusted prices} depend on the total fee only. With a tick, quotes cannot move by a fraction of it, so
the split changes \omterm{def:mx:fees-rebates-and-venue-economics:adjusted}{fee-adjusted prices} and queue lengths; it also matters when some traders do not pay fees directly (retail flow, some brokers).
\iqlookfor{Continuous-price neutrality; the tick; who bears the fee.}
\end{solution}

\begin{solution}{iq:mx:fees-rebates-and-venue-economics:3}
Different flow reaches the venue (routing order, fee-aware takers hitting another venue first), or different quoters and order types. Condition on the
same trades: compare fills from the same aggressive sweeps, or randomise the venue of your own passive orders and compare mark-outs.
\iqlookfor{Routing order; selection; a randomised test.}
\end{solution}

\begin{solution}{iq:mx:fees-rebates-and-venue-economics:4}
They become the back of the consolidated queue: they fill only when the cheaper venues are exhausted, less often and against more informed flow, and
their rebate must compensate.
\iqlookfor{Queue order across venues; adverse selection.}
\end{solution}

\begin{solution}{iq:mx:fees-rebates-and-venue-economics:5}
Randomly assign stocks to fee regimes with a control group (as the Transaction Fee Pilot intended), measure spreads, depth, fill rates and mark-outs,
routing shares and volumes, and estimate \omterm{def:mx:tick-size-queues-and-priority:did}{difference-in-differences} with clustered standard errors; check spillovers across venues.
\iqlookfor{Randomisation; control; spillovers.}
\end{solution}

\begin{solution}{iq:mx:fees-rebates-and-venue-economics:6}
Rebates must fall below 0.10 cents, so the maker-taker venues' advantage for providers shrinks, queues there shorten, and the inverted venues' edge for
takers narrows; my P\&L loses the rebate income and must come from spread capture and queue position, so I requote where fills are less toxic.
\iqlookfor{Fee arithmetic; queue response; P\&L components.}
\end{solution}
